\section{Pythagorean Carrier, Scale Fibers, and Unit-Circle Projection}
\label{sec:unified-reality-pythagorean-common-carrier}

The common-carrier discipline becomes concrete in the elementary
Pythagorean surface. The equation
$$
a^2+b^2=c^2
$$
and the normalized unit-circle equation
$$
(a/c)^2+(b/c)^2=1
$$
are not two unrelated facts. They are two reality readouts of one
carrier: an algebraic cone readout, a normalized shape readout, and a
scale readout whose loss is precisely the fiber of the shape projection.

\begin{definition}[Field-like reality interface]
\label{def:unified-reality-field-like-interface}
A \emph{field-like reality interface} is a stable reality interface with
a carrier $K$, classifier $\sim_K$, operations
$$
+,\ -,\ \cdot,\ 0,\ 1,
$$
and an apartness predicate $x\#0$. Its ledger licenses an inverse
$x^{-1}$ only from an apartness witness $x\#0$, and division is the
derived operation
$$
a/x:=a\cdot x^{-1}.
$$
Thus division is not a total primitive operation; it consumes an
apartness row.
\end{definition}

\begin{definition}[Nonzero scale carrier]
\label{def:unified-reality-nonzero-scale-carrier}
The nonzero scale carrier of $K$ is
$$
K^\#:=\sum_{c:K}(c\#0).
$$
An element is written $(c,\nu_c)$, where $\nu_c:c\#0$ is the displayed
apartness witness.
\end{definition}

\begin{definition}[Pythagorean carrier]
\label{def:unified-reality-pythagorean-carrier}
The \emph{Pythagorean carrier} over $K$ is
$$
\mathsf{Pyth}_K
:=
\sum_{a,b,c:K}
\bigl((c\#0)\times(a^2+b^2\sim_K c^2)\bigr).
$$
An element is written
$$
P=(a,b,c;\nu_c;h_P),
$$
where $\nu_c:c\#0$ licenses the divisions by $c$, and
$h_P:a^2+b^2\sim_K c^2$ is the square-sum certificate. The carrier is
therefore a packet with a scale-apartness row and an equation row, not a
bare triple.
\end{definition}

\begin{definition}[Algebraic unit-circle carrier]
\label{def:unified-reality-unit-circle-carrier}
The algebraic unit-circle carrier over $K$ is
$$
S^1_K
:=
\sum_{x,y:K}
\bigl(x^2+y^2\sim_K1\bigr).
$$
An element is written $Q=(x,y;h_Q)$. This is an algebraic circle
carrier. Topology, metric structure, arc length, and measure enter only
through additional reality interfaces.
\end{definition}

\begin{definition}[Pythagorean projections]
\label{def:unified-reality-pythagorean-projections}
The Pythagorean carrier has the following paper-side projections:
$$
\pi_{\mathrm{alg}}(a,b,c;\nu_c;h_P):=(a,b,c;h_P),
$$
$$
\pi_{\mathrm{scale}}(a,b,c;\nu_c;h_P):=(c,\nu_c),
$$
and
$$
\pi_{\mathrm{shape}}(a,b,c;\nu_c;h_P)
:=
\bigl(a/c,\ b/c;\ h_{\mathrm{shape}}\bigr).
$$
The proof row $h_{\mathrm{shape}}$ is supplied by
\autoref{thm:unified-reality-pythagorean-normalization}.
\end{definition}

\begin{theorem}[Pythagorean normalization]
\label{thm:unified-reality-pythagorean-normalization}
Let
$$
P=(a,b,c;\nu_c;h_P):\mathsf{Pyth}_K.
$$
Set
$$
x:=a/c,\qquad y:=b/c.
$$
Then
$$
x^2+y^2\sim_K1.
$$
Consequently $\pi_{\mathrm{shape}}$ is a well-defined projection from
$\mathsf{Pyth}_K$ to $S^1_K$.
\end{theorem}

\begin{proof}
By definition of division, $x=a\cdot c^{-1}$ and
$y=b\cdot c^{-1}$. The field-like ledger gives the usual distributive
and inverse laws under $\sim_K$, so
$$
x^2+y^2
\sim_K
a^2c^{-2}+b^2c^{-2}
\sim_K
(a^2+b^2)c^{-2}.
$$
Using $h_P:a^2+b^2\sim_Kc^2$, this becomes
$$
(a^2+b^2)c^{-2}
\sim_K
c^2c^{-2}
\sim_K
1,
$$
where the final step uses the apartness witness $\nu_c$.
\end{proof}

\concretizedIn{ch:concrete-instances-pythagorean-common-carrier-namecert}{Realizes the Pythagorean cone, nonzero scale, and unit-circle shape rows as one NameCert carrier.}

\begin{definition}[Scale reconstruction]
\label{def:unified-reality-pythagorean-scale-reconstruction}
Let $Q=(x,y;h_Q):S^1_K$ and let $(c,\nu_c):K^\#$. The
\emph{scale reconstruction} is
$$
\mathsf{Recon}(Q,c)
:=
(cx,cy,c;\nu_c;h_{\mathrm{recon}})
\in\mathsf{Pyth}_K,
$$
where $h_{\mathrm{recon}}$ is the square-sum certificate
$$
(cx)^2+(cy)^2\sim_Kc^2.
$$
\end{definition}

\begin{theorem}[Scale reconstruction]
\label{thm:unified-reality-pythagorean-scale-reconstruction}
For every $Q=(x,y;h_Q):S^1_K$ and $(c,\nu_c):K^\#$,
$$
(cx)^2+(cy)^2\sim_Kc^2.
$$
\end{theorem}

\begin{proof}
From $h_Q$ we have $x^2+y^2\sim_K1$. Therefore
$$
(cx)^2+(cy)^2
\sim_K
c^2x^2+c^2y^2
\sim_K
c^2(x^2+y^2)
\sim_K
c^2\cdot1
\sim_K
c^2.
$$
\end{proof}

\begin{theorem}[Shape-scale decomposition]
\label{thm:unified-reality-pythagorean-shape-scale-decomposition}
There is a classifier-respecting equivalence
$$
\mathsf{Pyth}_K\simeq S^1_K\times K^\#.
$$
The forward map is
$$
\Phi(P):=(\pi_{\mathrm{shape}}(P),\pi_{\mathrm{scale}}(P)),
$$
and the reverse map is
$$
\Psi(Q,c):=\mathsf{Recon}(Q,c).
$$
\end{theorem}

\begin{proof}
Let $P=(a,b,c;\nu_c;h_P)$. Applying $\Phi$ gives the shape-scale pair
$$
\left((a/c,b/c),\ (c,\nu_c)\right).
$$
Applying $\Psi$ returns
$$
(c(a/c),c(b/c),c;\nu_c;-).
$$
The inverse laws licensed by $\nu_c$ identify $c(a/c)\sim_Ka$ and
$c(b/c)\sim_Kb$, and the square-sum row transports along those
classifier witnesses. Hence $\Psi(\Phi(P))\sim_{\mathsf{Pyth}}P$.

Conversely, for $Q=(x,y;h_Q)$ and $(c,\nu_c):K^\#$, applying $\Psi$
gives $(cx,cy,c;\nu_c;-)$. Applying $\Phi$ returns
$$
\left(((cx)/c,(cy)/c;-),\ (c,\nu_c)\right).
$$
The same inverse laws identify $(cx)/c\sim_Kx$ and $(cy)/c\sim_Ky$, so
the shape row and scale row are classified together with $(Q,c)$.
\end{proof}

\begin{corollary}[Unit-circle readout is not the complete carrier]
\label{cor:unified-reality-unit-circle-readout-not-complete}
The shape carrier $S^1_K$ is not the complete Pythagorean carrier. The
complete carrier is classified by the shape together with the nonzero
scale:
$$
\mathsf{Pyth}_K\simeq S^1_K\times K^\#.
$$
\end{corollary}

\begin{proof}
This is exactly the decomposition of
\autoref{thm:unified-reality-pythagorean-shape-scale-decomposition}.
The projection to $S^1_K$ forgets the $K^\#$ coordinate.
\end{proof}

\begin{definition}[Scale fiber of a shape]
\label{def:unified-reality-pythagorean-scale-fiber}
For $Q:S^1_K$, the \emph{scale fiber} of the shape projection is
$$
\mathsf{ScaleFib}(Q)
:=
\sum_{P:\mathsf{Pyth}_K}
\bigl(\pi_{\mathrm{shape}}(P)\sim_{S^1_K}Q\bigr).
$$
\end{definition}

\begin{theorem}[Scale fiber classification]
\label{thm:unified-reality-pythagorean-scale-fiber-classification}
For every $Q:S^1_K$,
$$
\mathsf{ScaleFib}(Q)\simeq K^\#.
$$
\end{theorem}

\begin{proof}
Map $(P,\epsilon):\mathsf{ScaleFib}(Q)$ to
$\pi_{\mathrm{scale}}(P):K^\#$. Conversely, map $(c,\nu_c):K^\#$ to
the reconstructed packet $\mathsf{Recon}(Q,c)$ together with the
shape-classifier witness supplied by the inverse laws in the proof of
\autoref{thm:unified-reality-pythagorean-shape-scale-decomposition}.
The two composites are classified as identities by that same
shape-scale decomposition.
\end{proof}

\begin{corollary}[Shape projection is lossy when scales differ]
\label{cor:unified-reality-pythagorean-shape-projection-lossy}
If $Q:S^1_K$ is a displayed shape and $K^\#$ contains two scales $c,d$
with $c\not\sim_K d$, then the fiber of
$\pi_{\mathrm{shape}}:\mathsf{Pyth}_K\to S^1_K$ over $Q$ is
nontrivial. Consequently the shape projection satisfies the noncollapse
pattern of \autoref{def:unified-reality-projection-fiber} at every such
displayed shape.
\end{corollary}

\begin{proof}
By \autoref{thm:unified-reality-pythagorean-scale-fiber-classification},
the fiber over the displayed $Q$ is classified by $K^\#$. Distinct
scales therefore give distinct Pythagorean packets with the same shape
readout.
\end{proof}

\begin{definition}[Pythagorean extension lift]
\label{def:unified-reality-pythagorean-extension-lift}
Let $F\subset E$ be a field-extension reality with a displayed
extension ledger. If $a,b\in F$ and
$$
c\in E,\qquad c^2\sim_Ea^2+b^2,\qquad c\notin F,
$$
then $(a,b,c)$ is a \emph{Pythagorean extension lift} of the
$F$-data into $E$.
\end{definition}

\begin{theorem}[Extension lift reaches normalized coordinates]
\label{thm:unified-reality-pythagorean-extension-lift-shape}
Let $F\subset E$ be a field-extension reality. Suppose
$a,b\in F$ and
$$
c\in E,\qquad c\#0,\qquad c^2\sim_Ea^2+b^2,\qquad c\notin F.
$$
If $a\#0$, then $a/c\notin F$. If $b\#0$, then $b/c\notin F$.
Thus, in the nondegenerate case $a\#0$ and $b\#0$, the normalized shape
point $(a/c,b/c)$ belongs to the extension-facing carrier $E^2$ and not
to $F^2$.
\end{theorem}

\begin{proof}
Assume $a\#0$ and suppose $a/c\in F$. Write $r:=a/c$. Since $a\#0$
and $c\#0$, the field-like ledger gives $r\#0$. Then
$$
c\sim_E a/r.
$$
Because $a\in F$, $r\in F$, and $r\#0$, closure of the field-like
interface $F$ under licensed division gives $a/r\in F$, hence
$c\in F$, contradicting the extension-lift hypothesis. The argument for
$b/c$ is identical.
\end{proof}

\begin{example}[The \texorpdfstring{$(1,1,\sqrt{2})$}{(1,1,sqrt2)} carrier]
\label{ex:unified-reality-pythagorean-sqrt-two}
Let $F=\mathbb Q$ and $E=\mathbb Q(\sqrt{2})$. The packet
$$
(1,1,\sqrt{2})
$$
satisfies $1^2+1^2=(\sqrt{2})^2$. Its normalized shape is
$$
\left(1/\sqrt{2},1/\sqrt{2}\right).
$$
Since $1/\sqrt{2}=\sqrt{2}/2\notin\mathbb Q$, the scale readout and the
unit-circle readout both consume the same extension ledger
$\mathbb Q\subset\mathbb Q(\sqrt{2})$.
\end{example}

\begin{definition}[Pythagorean measure readouts]
\label{def:unified-reality-pythagorean-measure-readouts}
The triangle-area readout is the algebraic projection
$$
A_{\triangle}(a,b,c;\nu_c;h_P):=\frac{1}{2}ab.
$$
It is an oriented algebraic area row unless positivity and orientation
ledgers are added. If $K$ is embedded in an analytic reality $A$ whose
ledger supplies a circle constant, the disk-area readout is
$$
A_{\circ}(a,b,c;\nu_c;h_P):=\pi c^2
$$
inside $A$. The constant $\pi$ is not supplied by the field-like
interface $K$; it belongs to the analytic reality that certifies it.
\end{definition}

\begin{theorem}[Area readout through shape and scale]
\label{thm:unified-reality-pythagorean-area-shape-scale}
Let $P=(a,b,c;\nu_c;h_P):\mathsf{Pyth}_K$, and let
$x=a/c$ and $y=b/c$. Then
$$
\frac{1}{2}ab\sim_K\frac{1}{2}c^2xy.
$$
\end{theorem}

\begin{proof}
The inverse laws give $a\sim_Kcx$ and $b\sim_Kcy$. Therefore
$$
ab\sim_K(cx)(cy)\sim_Kc^2xy.
$$
Multiplying by $\frac{1}{2}$ gives the stated area readout.
\end{proof}

\begin{corollary}[Measure may require an analytic layer]
\label{cor:unified-reality-pythagorean-measure-analytic-layer}
Even when the scale belongs to an algebraic extension $E$, the disk-area
readout $\pi c^2$ may require an analytic reality strictly beyond $E$.
\end{corollary}

\begin{proof}
The algebraic extension ledger accounts for $c$ and hence for $c^2$.
It does not by itself account for the circle constant $\pi$. If the
active analytic reality supplies $\pi$ only through a separate process
or transcendental ledger, the disk-area readout factors through that
higher interface.
\end{proof}

\begin{definition}[Pythagorean generalized continuation]
\label{def:unified-reality-pythagorean-generalized-continuation}
The Pythagorean algebra-to-shape continuation packet is
$$
\mathsf{GenCont}_{\mathrm{alg}\to\mathrm{shape}}
:=
(\mathsf{Pyth}_K,\pi_{\mathrm{alg}},\pi_{\mathrm{shape}},
\mathsf{Tr}_{\mathrm{alg},\mathrm{shape}},
\mathsf{Led}_{\mathrm{alg},\mathrm{shape}}),
$$
where the transport row is
\autoref{thm:unified-reality-pythagorean-normalization}. The reverse
shape-scale-to-algebra packet uses
\autoref{thm:unified-reality-pythagorean-scale-reconstruction}.
\end{definition}

\begin{theorem}[Pythagorean algebra-geometry transport]
\label{thm:unified-reality-pythagorean-algebra-geometry-transport}
On the common carrier $\mathsf{Pyth}_K$, the algebraic equation
$a^2+b^2\sim_Kc^2$ and the normalized shape equation
$(a/c)^2+(b/c)^2\sim_K1$ are transport-equivalent only with the scale
fiber retained. The unit-circle point alone does not reconstruct the
Pythagorean packet.
\end{theorem}

\begin{proof}
The algebra-to-shape direction is
\autoref{thm:unified-reality-pythagorean-normalization}. The reverse
direction is
\autoref{thm:unified-reality-pythagorean-scale-reconstruction}, whose
input is not a bare point of $S^1_K$ but a pair in $S^1_K\times K^\#$.
By \autoref{cor:unified-reality-unit-circle-readout-not-complete}, the
scale fiber is required for the full carrier read.
\end{proof}

\begin{theorem}[Pythagorean common-carrier principle]
\label{thm:unified-reality-pythagorean-common-carrier-principle}
The carrier $\mathsf{Pyth}_K$ binds algebraic equation, normalized
shape, nonzero scale, triangle area, and disk area as projections of one
packet. It satisfies
$$
\mathsf{Pyth}_K\simeq S^1_K\times K^\#,
$$
and the projection to $S^1_K$ has scale fiber $K^\#$. If the scale is
read only through a field extension, the normalized shape coordinates
consume the same extension ledger in the nondegenerate case; if disk
area is read, a further analytic ledger may be required.
\end{theorem}

\begin{proof}
The equivalence is
\autoref{thm:unified-reality-pythagorean-shape-scale-decomposition}, and
the fiber computation is
\autoref{thm:unified-reality-pythagorean-scale-fiber-classification}.
The field-extension statement is
\autoref{thm:unified-reality-pythagorean-extension-lift-shape}. The
measure statement is
\autoref{cor:unified-reality-pythagorean-measure-analytic-layer}. All
five readouts use the same source row
$(a,b,c;\nu_c;h_P)$, so they are reality projections of one common
carrier rather than independent objects.
\end{proof}
